% 385_Higgs_Vakuum_Weg_De_ch.tex
% Johann Pascher, 1. Okt. 2026
% Der Higgs-Vakuum-Weg zu xi: Herkunft, Formeln, Status

\providecommand{\BM}{\textbf{[B]}}
\providecommand{\KM}{\textbf{[K]}}
\providecommand{\SM}{\textbf{[S]}}
\providecommand{\XM}{\textbf{[X]}}
\providecommand{\QM}{\textbf{[Q]}}

\chapter*{Der Higgs-Vakuum-Weg zu \texorpdfstring{$\xi$}{xi}}
\addcontentsline{toc}{chapter}{Der Higgs-Vakuum-Weg zu xi}

\begin{abstract}
Der erste Weg zu den Zusammenhängen der FFGFT führte über das Vakuum und das Higgs-Feld. Im
Frühjahr 2025 entstand daraus eine Formel, die die Higgs-Größen mit der elektrischen
Feldkonstante $\varepsilon_0$ verband:
$\lambda_h^2v^2e^2/(64\pi^4\varepsilon_0\hbar c\,m_h^2)$. Dieses Dokument verfolgt die Formel
anhand der Versionsgeschichte vom Februar 2025 bis heute, rechnet jede Stufe nach und stellt den
heutigen Stand fest. Ausgeschrieben ist die Vakuumformel genau $\xi_{\text{EFT}}\cdot\alpha$ mit
$\xi_{\text{EFT}}=\lambda_h^2v^2/(16\pi^3m_h^2)$; mit der FFGFT-Konvention $\alpha=1$ wird sie zum
heutigen Higgs-Weg. Mit der Selbstkopplung $\lambda_h=m_h^2/(2v^2)$ ergibt sich
$\xi_{\text{EFT}}=m_h^2/(64\pi^3v^2)=1{,}30\times10^{-4}$, rund $2{,}3\,\%$ unter $\xi=4/30000$.
Die Higgs-Geometrie ist damit eine Konsistenzprüfung von $\xi$ auf wenige Prozent \KM, keine
exakte Herleitung. Nicht tragfähig sind die frühe Gleichsetzung des Ausdrucks mit~1, eine kurzlebige
Variante mit $\mu_0/\varepsilon_0$ und die spätere Verkürzung auf $1/(16\pi^3)$ \XM. Der Faktor
$64\pi^4$ der frühen HTML-Fassungen stammt aus dem $4\pi$ in $4\pi\varepsilon_0$.
\end{abstract}

% ============================================================
\section*{Statusmarkierungen}
% ============================================================
\begin{center}
\begin{tabular}{@{}lp{11cm}@{}}
\textbf{[B]} & Algebraisch bewiesen. \\[2pt]
\textbf{[K]} & Numerisch bestätigt (Prüfskript, Vergleich mit Messwerten). \\[2pt]
\textbf{[S]} & Setzung, offen oder spekulativ. \\[2pt]
\textbf{[X]} & Widerlegt oder nicht tragfähig. \\[2pt]
\textbf{[Q]} & Konvention. \\
\end{tabular}
\end{center}
Vergleichswerte: CODATA 2022, PDG 2024 ($m_h=125{,}20$\,GeV, $v=246{,}22$\,GeV). Fundstellen der
frühen Formeln sind Commits des Repositorys; die Dateien existieren heute nur noch in der
Versionsgeschichte. Prüfskript: \texttt{pruef\_385\_higgs\_vakuum.py}.

% ============================================================
\section{Ausgangsidee: Vakuum, Higgs-Feld und Lichtgeschwindigkeit}
% ============================================================

Im Februar 2025 stand am Anfang eine qualitative These: Das Vakuum ist nicht leer, und das
Higgs-Feld zusammen mit der Hintergrundstrahlung bedingt die maximale Geschwindigkeit. Die
Verbindung lief über die Vakuumkonstanten,
\begin{equation}
  c=\frac{1}{\sqrt{\mu_0\varepsilon_0}},\qquad Z_0=\sqrt{\frac{\mu_0}{\varepsilon_0}}=376{,}73\ \Omega ,
\end{equation}
noch ohne Formel für einen Parameter (Erst-Upload vom 27.~März 2025, Dokument zur
Schwarzschild-Metrik). Parallel entstanden Notizen, die $\mu_0$ und $\varepsilon_0$ aus $\alpha$,
$e$ und $c$ ausdrücken, $\alpha=e^2/(4\pi\varepsilon_0\hbar c)$. Diese beiden Linien, Higgs und
Vakuum auf der einen, Vakuumkonstanten und $\alpha$ auf der anderen Seite, sind der Ursprung des
Weges, der hier verfolgt wird.

% ============================================================
\section{Die Formelkette 2025}
% ============================================================

\begin{center}
\small
\begin{tabular}{@{}p{2.2cm}p{6.4cm}p{5.4cm}@{}}
\toprule
\textbf{Datum} & \textbf{Formel} & \textbf{Befund} \\
\midrule
31.\,3.\,2025 & $\beta=\dfrac{\lambda_h^2v^2}{16\pi^3c^3}\cdot\dfrac{\hbar}{m_h^2}\cdot\dfrac1{r_0}$ & Faktor $16\pi^3$ entsteht; in SI nicht dimensionslos \\[8pt]
5.\,4.\,2025 & $\xi=\dfrac{\lambda_h^2v^2}{16\pi^3m_h^2}\approx1{,}33\times10^{-4}$ & richtig ausgewertet $1{,}30\times10^{-4}$ \\[8pt]
6.\,4.\,2025 & $\dfrac{\lambda_h^2v^2e^2}{64\pi^4\varepsilon_0\hbar c\,m_h^2}=1$ & $=\xi_{\text{EFT}}\,\alpha$, nicht 1 \\[8pt]
18.\,4.\,2025 & $\beta_T=\dfrac{\lambda_h^2v^2}{4\pi^2\lambda_0^2\alpha_0}$, Dimension $(\mu_0/\varepsilon_0)^k$ & dimensional nicht stimmig \\[8pt]
25.\,5.\,2025 & $\varepsilon_0$ über $e^2=4\pi\varepsilon_0$ eliminiert: $64\pi^4\to16\pi^3$ & Rückführung auf den Higgs-Ausdruck \\[8pt]
ab 1.\,6.\,2025 & $\xi=\dfrac{\lambda_h^2v^2}{64\pi^4m_h^2}=1{,}04\times10^{-5}$ & $4\pi$ aus $4\pi\varepsilon_0$ stehen geblieben \\[8pt]
22.\,8.\,2026 & $\lambda_h=m_h/v\Rightarrow\xi=1/(16\pi^3)$, Iteration zu $4/30000$ & Verkürzung, nicht tragfähig \\
\bottomrule
\end{tabular}
\end{center}

\subsection{Der Higgs-Ausdruck}

Am 31.~März 2025 erscheint zum ersten Mal ein Ausdruck mit den Higgs-Größen und dem Faktor
$16\pi^3$, damals als Parameter $\beta$ mit einem Radius $r_0$. Am 5.~April 2025 wird daraus mit
$r_0=\xi\,\ell_P$, natürlichen Einheiten und $\beta_T=1$
\begin{equation}
  \xi=\frac{\lambda_h^2v^2}{16\pi^3m_h^2}.
\end{equation}
Das Dokument nennt $1{,}33\times10^{-4}$; mit den Eingaben ausgerechnet sind es
$1{,}30\times10^{-4}$. Im selben Dokument steht bereits die Umformung mit $m_h^2=2\lambda_hv^2$:
$\xi=\lambda_h/(32\pi^3)$. Der Ausdruck ist dimensionslos und richtig gebildet \BM.

\subsection{Die Vakuumformel}

Am 6.~April 2025 verbindet das Dokument über natürliche Einheiten mit $\alpha=1$ den
Higgs-Ausdruck mit der Feinstrukturkonstante in ihrer SI-Form:
\begin{equation}
  \beta_T\,\alpha_{\text{EM}}\approx\frac{\xi\,\lambda_h^2v^2}{16\pi^3m_h^2}\cdot\frac1\xi\cdot
  \frac{e^2}{4\pi\varepsilon_0\hbar c}
  =\frac{\lambda_h^2v^2e^2}{64\pi^4\varepsilon_0\hbar c\,m_h^2}=1 .
  \label{eq:vakuum}
\end{equation}
Im selben Dokument stehen $c=1/\sqrt{\mu_0\varepsilon_0}$, $Z_0=\sqrt{\mu_0/\varepsilon_0}$ und
$h=2\pi\hbar\mu_0c^2$. Das ist die Formel, in der Higgs-Sektor und Vakuumkonstanten zum ersten Mal
zusammen auftreten.

\subsection{Die Variante mit \texorpdfstring{$\mu_0/\varepsilon_0$}{mu0/eps0}}

Am 18.~April 2025 versucht eine englische Arbeitsfassung eine \glqq elektromagnetische
Herleitung\grqq{} von $\beta_T$ aus $\mu_0$ und $\varepsilon_0$:
$\beta_T^{\text{nat}}=\lambda_h^2v^2/(4\pi^2\lambda_0^2\alpha_0)$ mit der Higgs-Compton-Länge
$\lambda_0=\hbar/(m_hc)$ und einem Dimensionsfaktor $(\mu_0/\varepsilon_0)^k$, $k=0$. Die Datei wurde
zwei Tage später gelöscht und nie übernommen.

\subsection{Rückführung und Ableger}

Am 25.~Mai 2025 wird $\varepsilon_0$ über $e^2=4\pi\varepsilon_0$ ($\alpha=1$, $\hbar=c=1$)
eliminiert; aus $64\pi^4$ wird wieder $16\pi^3$, und der Ausdruck ist wieder der Higgs-Ausdruck
von 5.~April. Ab 1.~Juni 2025 erscheint in den HTML-Fassungen und in Dok.~049 dagegen
$\xi=\lambda_h^2v^2/(64\pi^4m_h^2)=1{,}04\times10^{-5}$: Hier wurde $\varepsilon_0$ gestrichen,
das $4\pi$ aus $4\pi\varepsilon_0$ aber nicht gekürzt.

\subsection{Die Verkürzung von 2026}

Beim Anlegen von Dok.~320 am 22.~August 2026 wird $\lambda_h=m_h/v$ eingesetzt. Damit kürzt sich
alles auf $1/(16\pi^3)$, und eine Iteration soll von dort zum Fixpunkt $4/30000$ führen; als
Quellen sind A130 und Dok.~133 genannt. Diese Einsetzung verwechselt die Higgs-Selbstkopplung mit
einer Yukawa-Kopplung $y=m/v$; dieselbe Gleichsetzung $y=\lambda_h$ steht schon 2025 in
Dok.~097. A130 enthält den Higgs-Weg in keiner Fassung, Dok.~133 keine Iteration.

% ============================================================
\section{Was die Vakuumformel bedeutet}
% ============================================================

Ausgeschrieben ist der rechte Teil von \eqref{eq:vakuum} ein Produkt zweier dimensionsloser
Faktoren:
\begin{equation}
  \frac{\lambda_h^2v^2e^2}{64\pi^4\varepsilon_0\hbar c\,m_h^2}
  =\underbrace{\frac{\lambda_h^2v^2}{16\pi^3m_h^2}}_{\xi_{\text{EFT}}}\cdot
   \underbrace{\frac{e^2}{4\pi\varepsilon_0\hbar c}}_{\alpha}
  =\xi_{\text{EFT}}\cdot\alpha\quad\BM .
\end{equation}
Die Formel ist dimensional korrekt. Die Vakuumkonstante $\varepsilon_0$ tritt nur über $\alpha$
auf, und der Faktor $64\pi^4=16\pi^3\cdot4\pi$ entsteht, weil das $4\pi$ aus $4\pi\varepsilon_0$
in den Nenner wandert. In SI ergibt der Ausdruck $\xi_{\text{EFT}}\,\alpha=9{,}5\times10^{-7}$
\KM. Mit der FFGFT-Konvention $\alpha=1$, also $e^2=4\pi\varepsilon_0\hbar c$, ist er genau
$\xi_{\text{EFT}}$ \BM. Der Higgs-Vakuum-Weg und der heutige Higgs-Weg sind damit dieselbe
Rechnung: Die Vakuumkonstanten verschwinden in natürlichen Einheiten, übrig bleibt der
Higgs-Ausdruck, der $\xi$ prüft. So deutet es auch Dok.~310: In natürlichen Einheiten
verschwinden $\mu_0$ und $\varepsilon_0$, es bleibt nur $\xi$.

Nicht tragfähig ist die Gleichsetzung mit~1 in \eqref{eq:vakuum} \XM. Im ersten Schritt wird
$\xi\cdot(\ldots)\cdot(1/\xi)$ gekürzt, während $\beta_T$ selbst den Faktor $1/\xi$ enthält.
Richtig gelesen ist $\beta_T=\xi_{\text{EFT}}/\xi=0{,}977$, also fast, aber nicht genau~1, und
$\beta_T\alpha=0{,}0071$.

Die Variante mit $\mu_0/\varepsilon_0$ vom 18.~April ist dimensional nicht stimmig \XM:
$\lambda_0^{-2}$ hat die Dimension Energie$^2$, $v^2/\lambda_0^2$ also Energie$^4$; mit
$\lambda_0=1/m_h$ ergibt der Ausdruck rund $4\times10^5$\,GeV$^4$, weder~1 noch $\xi$. Der Faktor
$(\mu_0/\varepsilon_0)^k=Z_0^{2k}$ steht nur in der Dimensionszeile und verschwindet mit $k=0$.

% ============================================================
\section{Heutiger Stand}
% ============================================================

Mit der Selbstkopplung des Standardmodells, $\lambda_h=m_h^2/(2v^2)=0{,}129$, wird der
Higgs-Ausdruck zu
\begin{equation}
  \xi_{\text{EFT}}=\frac{\lambda_h^2v^2}{16\pi^3m_h^2}=\frac{m_h^2}{64\pi^3v^2}
  =\frac{\lambda_h}{32\pi^3}=1{,}30\times10^{-4}\quad\KM ,
\end{equation}
rund $2{,}3\,\%$ unter $\xi=4/30000=1{,}333\times10^{-4}$. Er hängt nur vom Verhältnis $m_h/v$ ab.
Je nach Eingaben liegt der Abstand zwischen $2{,}2$ und $2{,}4\,\%$; die oft genannten $2{,}6\,\%$
entstehen nur mit dem auf $0{,}129$ gerundeten $\lambda_h$:
\begin{center}
\small
\begin{tabular}{@{}lcc@{}}
\toprule
\textbf{Eingaben} & $\xi_{\text{EFT}}$ & \textbf{Abstand zu }$4/30000$ \\
\midrule
$m_h=125{,}20$, $v=246{,}22$\,GeV (PDG 2024) & $1{,}303\times10^{-4}$ & $-2{,}3\,\%$ \\
$m_h=125{,}25$, $v=246{,}22$\,GeV & $1{,}304\times10^{-4}$ & $-2{,}2\,\%$ \\
$m_h=125{,}1$, $v=246{,}2$\,GeV & $1{,}301\times10^{-4}$ & $-2{,}4\,\%$ \\
$\lambda_h=0{,}129$ gerundet, $m_h=125{,}1$ & $1{,}299\times10^{-4}$ & $-2{,}6\,\%$ \\
\bottomrule
\end{tabular}
\end{center}
Die Higgs-Geometrie ist damit eine Konsistenzprüfung von $\xi$ auf wenige Prozent \KM\
(Dok.~354, A190, 310). Die Abweichung ist strukturell gedeutet \SM\ (R59): Die SI-Werte von $m_h$
und $v$ sind schemaabhängig, und fraktale Korrekturen lassen sich nicht sauber von der
Schleifenstruktur des Standardmodells trennen. Eine exakte Herleitung von $4/30000$ ist der
Higgs-Weg nicht. Er ist aber eine eigenständige Prüfung: Er benutzt nur die Higgs-Masse und die
elektroschwache Skala, keine Leptonmasse und keine Galois-Zählung.

% ============================================================
\section{Was trägt und was nicht}
% ============================================================

\begin{center}
\small
\begin{tabular}{@{}p{7.2cm}p{1.2cm}p{5.4cm}@{}}
\toprule
\textbf{Aussage} & \textbf{Status} & \textbf{Grund} \\
\midrule
Vakuumformel $=\xi_{\text{EFT}}\cdot\alpha$ & \BM & Ausschreiben; $\alpha=e^2/(4\pi\varepsilon_0\hbar c)$ \\
mit $\alpha=1$ gleich dem Higgs-Ausdruck & \BM & $e^2=4\pi\varepsilon_0\hbar c$ \\
$\xi_{\text{EFT}}=m_h^2/(64\pi^3v^2)=1{,}30\times10^{-4}$, rund $2{,}3\,\%$ unter $\xi$ & \KM & Konsistenzprüfung \\
Abweichung strukturell (Schema, fraktale Korrektur) & \SM & R59 \\
Vakuumformel $=1$ & \XM & Kürzung von $1/\xi$ inkonsistent \\
Variante mit $\mu_0/\varepsilon_0$ & \XM & Dimension Energie$^4$ \\
$\xi=\lambda_h^2v^2/(64\pi^4m_h^2)=1{,}04\times10^{-5}$ & \XM & $4\pi$ aus $4\pi\varepsilon_0$ nicht gekürzt \\
$\lambda_h=m_h/v$, $\xi=1/(16\pi^3)$, Iteration zu $4/30000$ & \XM & 15-mal $\xi$; Iteration läuft gegen 0 \\
Higgs-Weg als exakte Herleitung von $4/30000$ & \XM & Abstand 2--3\,\% \\
\bottomrule
\end{tabular}
\end{center}

% ============================================================
\section{Fazit}
% ============================================================

Der Higgs-Vakuum-Weg, der am Anfang der FFGFT stand, trägt. Seine Vakuumformel von April 2025 ist,
richtig ausgeschrieben, das Produkt aus Higgs-Ausdruck und Feinstrukturkonstante; mit $\alpha=1$
ist sie der heutige Higgs-Weg, und dieser trifft $\xi$ auf rund $2{,}3\,\%$. Die Fehler der
Entwicklungslinie liegen in Nebenschritten: der Gleichsetzung mit~1, einer dimensional falschen
Variante, einem nicht gekürzten $4\pi$ und einer späten Verkürzung. Als Herleitung der exakten
Zahl $4/30000$ reicht der Weg nicht; als unabhängige Prüfung von $\xi$ aus dem Higgs-Sektor bleibt
er bestehen.

\vspace{1em}
\noindent\textit{Prüfskript:} \texttt{2/python/Dok385\_Skripte/pruef\_385\_higgs\_vakuum.py}
--- 17/17.
